\section{Introduction}
\label{ch:introduction}

% ============================================================================
% === WRITING GUIDE: Chapter 1, Introduction (~4 pages, ~1,800 words) =======
% ============================================================================
% JOB: establish the "missing E2R guideline" motivation, state the central
% thesis, present RQ1-RQ3 with hypotheses (H4 flagged untested up front), and
% list contributions. This chapter carries Mari Carmen's deferred review items
% "storyline framing" and "FACILE positioning" first.
%
% CENTRAL THESIS STATEMENT (recurring thread; open with a version of this):
%   When LLMs adapt figurative language into Easy-to-Read English, reliable
%   quality does not come from where current practice looks for it (larger
%   models, cleverer prompts) but from explicit structure: decomposing the
%   task into inspectable steps and grounding evaluation in human judgement.
%   The lever that pays is structure, in both the system and its measurement.
%
% WRITE THIS CHAPTER AFTER Discussion 8.1 (per WRITING_PLAN.md schedule),
% so the framing cannot overstate the evidence.
% ============================================================================

\subsection{Motivation}
% --- GUIDE (~1.25 pp): figurative language as an accessibility barrier.
% Must land:
%  (1) E2R guidelines explicitly prohibit figurative language, but no system
%      or guideline operationalises HOW to adapt it. Cite the UPM
%      per-guideline-module pattern: \cite{Surez-Figueroa2023} (morphological
%      features), Diab 2024 (dialogues), \cite{Surez-Figueroa2024} (FACILE).
%      This thesis is the figurative-language module of that pattern.
%  (2) The two-literatures gap: figurative-language NLP detects but never
%      adapts for accessibility; E2R automates other guidelines but not this
%      one. Neither community addresses the intersection (developed in Ch 2).
%  (3) Why LLMs now: supervised detection systems exist but replacement has
%      no training data and no automatic metric; zero/one-shot LLMs are the
%      only viable approach for the generative half of the task.

% TODO(jelle): prose.

\subsection{Problem Statement}
% --- GUIDE (~0.75 pp): state the task (given an English sentence, detect
% idiomatic/metaphorical expressions and produce a literal, meaning-preserving,
% easier-to-read restatement) and then the thesis statement paragraph.
% NOTE: the old bullet "unified treatment of idioms and metaphors" is resolved
% as PARALLEL treatment under a common architecture: four separate LangGraph
% workflows, prompts never shared across phenomena. Say "parallel", not
% "unified" (the implementation deliberately bifurcates).

% TODO(jelle): prose.

\subsection{Research Questions and Hypotheses}
% --- GUIDE (~1 pp): RQ1-RQ3 verbatim from wiki/thesis/research-questions.md:
%  RQ1: To what extent can a zero-/one-shot agentic LLM pipeline accurately
%       detect and interpret idiomatic and metaphorical expressions in
%       context, compared to established supervised benchmarks?
%  RQ2: Does an explicit agentic decomposition (detection -> explanation ->
%       transformation) improve figurative-language handling compared to
%       monolithic prompt-based LLM approaches?
%  RQ3: How effectively can an agentic LLM system replace idiomatic and
%       metaphorical expressions with literal, meaning-preserving
%       alternatives, and what trade-offs arise between semantic fidelity
%       and readability?
% Hypotheses H1-H4. H4 (explanation quality correlates with replacement
% correctness) MUST be stated here and explicitly flagged as untested in this
% thesis, with a pointer to the proposed design in the Discussion (8.5).
% Observability is framed as a design property examined qualitatively in the
% Discussion, NOT as a fourth RQ.

% TODO(jelle): prose.

\subsection{Contributions}
% --- GUIDE (~0.5 pp): four concrete contributions:
%  (1) An agentic open-weights pipeline for figurative-language E2R
%      adaptation (FastAPI + LangGraph + vLLM; 200+ persisted, reproducible
%      runs), with a monolithic baseline as designed control.
%  (2) An adapted Direct Assessment human-evaluation methodology for
%      figurative-language replacement (Ch 4), with documented departures
%      from Graham 2013/2017.
%  (3) Empirical findings on what does and does not transfer: scale lifts
%      detection but not human-perceived replacement quality; prompt gains
%      invert across scales; decomposition buys meaning preservation.
%  (4) A by-product dataset: 108 human-written alternative replacements +
%      rubric ratings for 81 system outputs.

% TODO(jelle): prose.

\subsection{Document Structure}
% --- GUIDE (~0.5 pp): one sentence per chapter. Close by framing the three
% empirical chapters as testing three levers: scale (RQ1), structure (RQ2),
% and task difficulty / fidelity-readability trade-offs (RQ3).

% TODO(jelle): prose.

\newpage
